\section{Study 1: Release-Height Generalisation}
\label{sec:study1}

This study investigates how \mcpilot{} scales across release heights from ground level
to 2.0\,m.
Raising the release point increases flight time, shifts the reachable target zone
outward, and changes the speed-to-distance mapping that the GP and policy must learn.
All experiments use the NumPy ballistic simulator (\texttt{mc-pilot-elevated/}) with
no sensor noise; only the trio $(z_{\mathrm{release}},\, \ell_M,\, T)$ varies from
the baseline; all other parameters follow the baseline values of Config~A.

\subsection{Experimental Setup}
\label{sec:study1:setup}

Five configurations span the height range of practical robotic throwing platforms.
Table~\ref{tab:height_configs} summarises the geometry of each; the target
distribution is five uniformly spaced positions along the sagittal axis,
$P_x \in [0.75\,\ell_M,\, \ell_M]$.

\begin{table}[t]
  \centering
  \caption{Release-height configurations. $N_{\mathrm{exp}} = 5$ for A; varied for B--E.
           ``Strat'' = stratified exploration.}
  \label{tab:height_configs}
  \small
  \begin{tabular}{lcccl}
    \toprule
    Config & $z_{\mathrm{rel}}$ & $\ell_M$ & $T$ & Result \\
    \midrule
    A          & 0.5\,m & 1.75\,m & 0.70\,s & 5/5 hits \\
    B (seed=1) & 1.0\,m & 1.90\,m & 0.85\,s & 0/10 \textbf{FAIL} (random) \\
    B (seed=2) & 1.0\,m & 1.90\,m & 0.85\,s & 10/10 (random) \\
    B-Strat    & 1.0\,m & 1.90\,m & 0.85\,s & 10/10 (stratified) \\
    C-Strat    & 1.5\,m & 2.15\,m & 0.95\,s & 10/10 \\
    D-Strat    & 2.0\,m & 2.35\,m & 1.00\,s & 10/10 \\
    E-Strat5   & 0.0\,m & 1.10\,m & 0.55\,s & 10/10 ($\ell_s = 0.15$\,m) \\
    \bottomrule
  \end{tabular}
\end{table}

\subsection{Results: Configs B, C, D}
\label{sec:study1:results_bcd}

Table~\ref{tab:bcd_errors} reports per-trial landing errors for the three successful
elevated configurations under stratified exploration.

\begin{table}[t]
  \centering
  \caption{Per-trial landing errors (mm) for Configs B-Strat through D-Strat.
           All 10/10 trials hit (error $<$ 50\,mm). Mean over 10 trials.}
  \label{tab:bcd_errors}
  \small
  \begin{tabular}{lccc}
    \toprule
    Trial & B-Strat (1.0\,m) & C-Strat (1.5\,m) & D-Strat (2.0\,m) \\
    \midrule
    1  & 6  & 16 & 7  \\
    2  & 26 & 5  & 19 \\
    3  & 29 & 12 & 15 \\
    4  & 18 & 22 & 25 \\
    5  & 11 & 27 & 21 \\
    6  & 28 & 7  & 12 \\
    7  & 37 & 11 & 18 \\
    8  & 13 & 11 & 9  \\
    9  & 17 & 24 & 29 \\
    10 & 12 & 8  & 9  \\
    \midrule
    Mean & $\approx 20$\,mm & $\approx 14$\,mm & $\approx 16$\,mm \\
    \bottomrule
  \end{tabular}
\end{table}

Three observations are consistent across B, C, D.
First, all configurations converge: the cost for D-Strat falls from 0.827 to 0.024
within the first trial --- the fastest first-trial convergence in this study.
Second, mean landing error is roughly constant across heights (14--20\,mm), indicating
that the GP generalises the additional flight time as height increases without
systematic accuracy degradation.
Third, error variance does not grow with height: the worst-case error for D-Strat
(29\,mm) is comparable to C-Strat (27\,mm).

\subsection{Config E: Progressive Diagnosis of RBF Blindness}
\label{sec:study1:config_e}

Config E ($z_{\mathrm{rel}} = 0.0$\,m, target range 0.35\,m) is the corner case where
the required speed range (approximately $[2.80,\, 3.39]$\,m/s --- only 17\% of
$[0,\, u_M]$) is narrow relative to the default RBF lengthscale.
Five consecutive diagnosis iterations were required to isolate and fix the failure;
the progression illustrates the general diagnostic methodology.

\paragraph{E-Strat (0/10): coverage failure.}
Standard stratified exploration over $[0, u_M]$ fails: the cost plateaus near 0.25.
Only the highest-speed band produces throws near the targets; lower bands provide no
GP supervision for the relevant speed range.

\paragraph{E-Strat2 (0/10): cost plateau persists after targeted coverage.}
Restricting stratification to $[2.5,\, u_M]$ ensures all throws land near targets,
yet the cost plateau persists.
This rules out data coverage and implicates the policy architecture.

\paragraph{E-Strat3 (0/10): doubling exploration fails.}
$N_{\mathrm{exp}} = 10$ with targeted stratification leaves the plateau unchanged.
Policy inspection reveals constant output $u_M = 3.5$\,m/s for every target --- the
$\tanh$ saturates because all RBF centres produce the same activation in the
0.35\,m-wide zone at $\ell_s = 1.0$\,m.

\paragraph{E-Strat4 (0/10): tighter cost does not help.}
Reducing $\ell_c = 0.25$\,m to sharpen gradient signals leaves the policy outputting
constant speed for every trial.
The failure is unambiguously the RBF lengthscale.

\paragraph{E-Strat5 (10/10, mean 26\,mm): $\ell_s = 0.15$\,m resolves blindness.}
Setting $\ell_s = 0.15$\,m per Rule~\ref{eq:ls_rule} immediately achieves 10/10 hits
with errors in the range 20--33\,mm.
This is the tightest absolute error band in the study, consistent with higher spatial
resolution in the narrow target zone.
See Section~\ref{sec:convergence:lengthscale} for the derivation of the rule.

\subsection{Summary}
\label{sec:study1:summary}

After applying the two convergence rules (Section~\ref{sec:convergence}), \mcpilot{}
achieves 100\% hit rates on all five height configurations with mean errors below
35\,mm.
The key parameter adjustments are (i) scaling $\ell_M$ and $T$ geometrically with
height and (ii) scaling $\ell_s$ inversely with target range.
No per-configuration tuning of the GP kernel, particle count, or optimisation
schedule is required.
